\documentclass[11pt]{article}
\usepackage[T1]{fontenc}
\usepackage{amsmath,amssymb,amsthm}
\usepackage[margin=1in]{geometry}
\newcommand{\E}{\mathbb{E}}
\newcommand{\Prb}{\mathbb{P}}
\newcommand{\SAT}{\mathrm{SAT}}
\newcommand{\VNP}{\mathrm{VNP}}
\newcommand{\VP}{\mathrm{VP}}
\newcommand{\F}{\mathbb{F}}
\DeclareMathOperator{\Var}{Var}
\DeclareMathOperator{\Cov}{Cov}
\DeclareMathOperator{\Gal}{Gal}
\DeclareMathOperator{\Spl}{Spl}
\DeclareMathOperator{\size}{size}

\title{SAT Partition Functions and Algebraic Complexity:\\Formula Notebook}
\author{AI-assisted research notes; no claim of a P vs NP proof}
\date{24 November 2026}
\begin{document}
\maketitle

\noindent
\textbf{Status notation.}
[D] definition; [I] algebraic identity; [T] theorem under stated hypotheses;
[C] verified example/counterexample; [O] open research question.
No statement in this notebook proves $\mathrm{P}\ne\mathrm{NP}$ or $\VP\ne\VNP$.

\section{Boolean SAT and a partition polynomial [D, I]}
Let $\varphi=C_1\land\cdots\land C_m$ be a CNF formula
on $n$ Boolean variables, and let
\[
I_j(y)=
\begin{cases}
1,&C_j(y)=\mathrm{false},\\
0,&C_j(y)=\mathrm{true},
\end{cases}
\qquad
E_\varphi(y)=\sum_{j=1}^{m}I_j(y).
\]
Define
\[
Z_\varphi(z)=\sum_{y\in\{0,1\}^{n}}z^{E_\varphi(y)}
=\sum_{k=0}^{m}c_kz^k,\qquad
c_k=\#\{y:E_\varphi(y)=k\}.
\]
With $0^0=1$,
\[
Z_\varphi(0)=c_0=\#\SAT(\varphi),\qquad
Z_\varphi(1)=2^n,\qquad
\sum_{k=0}^{m}c_k=2^n.
\]
These are counting identities; they do not assert efficient
computation of $c_0$.

\section{Characteristic function, Fourier inversion, quantum amplitude [I]}
Let $Y$ be uniformly distributed on $\{0,1\}^n$. Define
\[
p_k=\Prb(E_\varphi(Y)=k)=2^{-n}c_k,
\qquad
\mathcal{A}_\varphi(\theta)
=\E[e^{i\theta E_\varphi(Y)}]
=2^{-n}Z_\varphi(e^{i\theta})
=\sum_{k=0}^{m}p_ke^{ik\theta}.
\]
For $\omega=e^{2\pi i/(m+1)}$,
\[
c_0=\frac{1}{m+1}\sum_{j=0}^{m}Z_\varphi(\omega^j)
=\frac{2^n}{m+1}\sum_{j=0}^{m}
\mathcal{A}_\varphi\!\left(\frac{2\pi j}{m+1}\right).
\]
If $U_\theta|y\rangle=e^{i\theta E_\varphi(y)}|y\rangle$,
then
\[
\langle 0^n|H^{\otimes n}U_\theta H^{\otimes n}|0^n\rangle
=\mathcal{A}_\varphi(\theta).
\]
This amplitude identity implies neither a quantum speedup nor a
complexity lower bound.

\section{Exact low-order statistics [I]}
For each fixed positive integer $r$,
\[
\E[E_\varphi(Y)^r]
=\sum_{j_1,\ldots,j_r=1}^{m}
\Prb(I_{j_1}(Y)=\cdots=I_{j_r}(Y)=1).
\]
For any specified tuple of clauses, the joint violation
probability is $0$ if the required literal assignments conflict,
and otherwise $2^{-t}$, where $t$ is the number of distinct
Boolean variables fixed by those requirements.
Consequently, for fixed $r$, the moment can be evaluated exactly
with $O(m^r\operatorname{poly}(n,m))$ elementary operations.

If every clause has exactly three distinct variables, then
\[
\E[E_\varphi(Y)]=\frac m8.
\]
In general,
\[
\Var(E_\varphi)
=\sum_j\Var(I_j)+2\sum_{j<k}\Cov(I_j,I_k).
\]
For the star instance
\[
\varphi_\star=\bigwedge_{j=1}^{m}(x\lor a_j\lor b_j),
\]
where all $a_j,b_j$ are distinct and different from $x$,
\[
\Var(E_{\varphi_\star})
=\frac{m^2+6m}{64}.
\]

\section{Cumulants and a local convergence guarantee [T]}
Let the clause dependency graph join clauses sharing a Boolean
variable, let $\Delta$ be its maximum degree and $D=\Delta+1$.
For $r\ge2$, the bounded-variable dependency-graph cumulant
bound of F\'eray--M\'eliot--Nikeghbali yields
\[
|\kappa_r(E_\varphi)|
\le m\,r^{r-2}(2D)^{r-1}
\le m\,r![2eD]^{r-1}.
\]
Writing $C=2eD$, for $|\theta|<C^{-1}$, the branch of
$\log\mathcal{A}_\varphi$ analytic at $\theta=0$ satisfies
\[
\log\mathcal{A}_\varphi(\theta)
=\sum_{r=1}^{\infty}
\frac{(i\theta)^r\kappa_r(E_\varphi)}{r!},
\]
and, for $R\ge1$,
\[
\left|
\log\mathcal{A}_\varphi(\theta)
-\sum_{r=1}^{R}
\frac{(i\theta)^r\kappa_r(E_\varphi)}{r!}
\right|
\le
\frac{m}{C}
\frac{(C|\theta|)^{R+1}}{1-C|\theta|}.
\]
These are \emph{sufficient} convergence guarantees,
not the exact Taylor radius.

\section{Complex zeros and counterexamples [C]}
Two clauses on five variables,
\[
\varphi_2=(x\lor a\lor b)\land(\neg x\lor c\lor d),
\]
give
\[
Z_{\varphi_2}(z)=8(3+z),\qquad
\mathcal{A}_{\varphi_2}(\theta)=\frac{3+e^{i\theta}}4.
\]
Hence the Taylor radius of $\log\mathcal{A}_{\varphi_2}$
about $0$ is
\[
\rho=\sqrt{\pi^2+(\ln3)^2}>\pi,
\qquad\Delta=1.
\]
For $k$ pairwise disjoint 3-clauses,
\[
Z_\varphi(z)=(7+z)^k,\qquad
\rho=\sqrt{\pi^2+(\ln7)^2}>\pi,\qquad\Delta=0.
\]
Conversely, with
\[
\begin{aligned}
\chi={}&(x\lor y\lor\neg z)
\land(x\lor\neg y\lor z)\\
&\land(\neg x\lor y\lor z)
\land(\neg x\lor\neg y\lor\neg z),
\end{aligned}
\]
we have
\[
Z_\chi(z)=4(1+z),\qquad
\mathcal{A}_\chi(\theta)=\frac{1+e^{i\theta}}2,\qquad
\mathcal{A}_\chi(\pi)=0,\qquad\Delta=3.
\]
Thus bounded $\Delta$ alone does not determine whether
$\log\mathcal{A}_\varphi$ is analytic at $\theta=\pi$.

\section{Shannon quantities and search [D, I]}
\[
H(E_\varphi)
=-\sum_{k:p_k>0}p_k\log_2p_k.
\]
If $c_0>0$, an independent uniform random assignment succeeds
with probability $c_0/2^n$, so its expected number of trials is
\[
\E[T_{\mathrm{uniform}}]=\frac{2^n}{c_0},
\qquad
-\log_2\Prb(E_\varphi(Y)=0)=n-\log_2c_0.
\]
These are statements about uniform independent sampling,
not lower bounds for arbitrary SAT algorithms.
For a separately specified joint distribution of problem
instances $X$, target assignments $Y$, and advice $D$,
\[
I(Y;D\mid X)=H(Y\mid X)-H(Y\mid X,D).
\]
Information gain alone does not imply an efficient search algorithm.

\section{Symbolic polynomial family and VNP [D, I]}
Fix a field $\F$. For $n\ge3$, let $\mathcal{C}_n$ be the
$8\binom n3$ distinct 3-literal clauses obtained by choosing
three different variables and an independent sign for each.
For $C\in\mathcal{C}_n$, let $X_C$ be an indeterminate.
For a Boolean literal $\ell$, write $v_\ell(y)=y_i$
for $\ell=x_i$ and $v_\ell(y)=1-y_i$ for $\ell=\neg x_i$.
If $C=(\ell_1\lor\ell_2\lor\ell_3)$, let
\[
I_C(y)=\prod_{t=1}^{3}(1-v_{\ell_t}(y)).
\]
Define a \emph{multivariate symbolic family}
\[
\boxed{
F_n(\mathbf X,z)
=\sum_{y\in\{0,1\}^n}
\prod_{C\in\mathcal C_n}
\bigl(1+X_C(z-1)I_C(y)\bigr).
}
\]
For any 3-CNF $\varphi$ built from \emph{distinct} clauses
in $\mathcal C_n$, let $X_C=1$ for clauses present in $\varphi$
and $X_C=0$ otherwise. Then
\[
F_n(\mathbf 1_\varphi,z)=Z_\varphi(z)
\quad\text{in }\F[z].
\]
If $\operatorname{char}\F=p>0$, integer coefficients are
understood modulo $p$.

Writing
\[
g_n(\mathbf X,z,\mathbf y)
=\prod_{C\in\mathcal C_n}
\bigl(1+X_C(z-1)I_C(\mathbf y)\bigr),
\]
the polynomials $g_n$ have polynomially bounded degree,
variable count and arithmetic circuit size, and
\[
F_n(\mathbf X,z)=\sum_{\mathbf y\in\{0,1\}^n}
g_n(\mathbf X,z,\mathbf y).
\]
Consequently, by the standard p-definability criterion,
\[
\boxed{(F_n)_{n\ge3}\in\VNP_{\F}.}
\]
This does \emph{not} establish VNP-completeness,
a superpolynomial circuit lower bound, or
$(F_n)\notin\VP_{\F}$.

\section{Fixed-instance circuit warning [I]}
For each \emph{fixed} instance $\varphi$,
\[
Z_\varphi(z)=c_0+c_1z+\cdots+c_mz^m
\]
has an arithmetic circuit of $O(m)$ gates by Horner's rule,
when its coefficients $c_k$ may be used as field constants.
Computing those constants uniformly is a different task.
Therefore a standard unrestricted circuit-size lower bound
cannot be inferred for this fixed univariate polynomial.

\section{Galois examples [I]}
For nonconstant $Z_\varphi\in\mathbb{Q}[z]$, put
\[
G_\varphi=\Gal(\Spl(Z_\varphi)/\mathbb Q).
\]
Repeated copies of the unit clause $(x)$ give
\[
Z_\varphi(z)=1+z^m,\qquad
\Spl(Z_\varphi)=\mathbb Q(\zeta_{2m}),\qquad
G_\varphi\cong(\mathbb Z/2m\mathbb Z)^\times.
\]
Pairwise disjoint 3-clauses give $(7+z)^k$,
whose splitting-field Galois group is trivial.
These counterexamples do not determine the Galois group
of a typical random SAT-generated polynomial.

\section{Open implications [O]}
\[
(F_n)\in\VNP_{\F}
\quad\not\Rightarrow\quad
(F_n)\notin\VP_{\F}.
\]
\[
\text{Empirical SAT hardness correlation}
\quad\not\Rightarrow\quad
\mathrm{P}\ne\mathrm{NP}.
\]
\[
\VP\ne\VNP
\quad\not\Rightarrow\quad
\mathrm{P}\ne\mathrm{NP}
\quad\text{by any implication established here}.
\]
The last line records the absence of a proved bridge
\emph{in these notes}; it is not a logical non-implication theorem.

\end{document}
